\chapter{Cushing's Syndrome}
\index{Cushing's Syndrome}
\label{sec:Cushing's Syndrome}

\section{Overview}


\begin{itemize}[noitemsep]
  \item Cushing's syndrome results from chronic exposure to excess cortisol, regardless of the cause, with exogenous glucocorticoid use (iatrogenic) being the most common cause overall. Endogenous causes are divided into ACTH-dependent (70--80\%): Cushing's disease (pituitary adenoma secreting ACTH, 70\%) and ectopic ACTH syndrome (small cell lung cancer, bronchial carcinoid)
  \item And ACTH-independent (20--30\%): adrenal adenoma, adrenal carcinoma, and bilateral adrenal abnormal increase in cells.
\end{itemize}
\section{Causes}
The underlying mechanisms involve complex interactions between genetic predisposition and environmental factors. Disruption of normal physiological processes leads to the characteristic features of the condition. Multiple contributing factors may act synergistically to trigger or worsen the condition.
\section{Risk Factors}


\begin{itemize}[noitemsep]
  \item Genetic predisposition and family history
  \item Advancing age
  \item Lifestyle factors (diet, exercise, smoking, alcohol)
  \item Environmental exposures
  \item Underlying medical conditions
  \item Medication use
\end{itemize}
\section{Signs and Symptoms}

\subsection{Common symptoms}
\begin{itemize}[noitemsep]
  \item You may experience central (truncal) obesity, moon facies, dorsocervical fat pad (``buffalo hump''), supraclavicular fat pads, thin extremities, purple striae ($>$1 cm wide), easy bruising, proximal muscle weakness, high blood pressure, hyperglycemia/diabetes, osteoporosis, menstrual irregularity, hirsutism, and psychiatric disturbances.
  \item Pain or discomfort in the affected area
  \end{itemize}

\subsection{Emergency warning signs}
\begin{itemize}[noitemsep]
  \item Severe or worsening symptoms of cushing's syndrome require immediate medical evaluation
\end{itemize}
\section{Progression and Stages}
The condition may progress through different stages of severity over time. Early stages often involve mild or subtle symptoms that may be overlooked. Moderate stages involve more noticeable symptoms affecting daily function. Advanced stages may involve significant impairment and complications.
\section{Complications}
\begin{itemize}[noitemsep]
  \item Disease progression leading to worsening symptoms
  \item Secondary infections or injuries
  \item Side effects of treatment
  \item Reduced quality of life
  \item Emotional and psychological impact
\end{itemize}
\section{Diagnosis}
Diagnosis typically involves confirmation of hypercortisolism (24-hour urinary free cortisol, late-night salivary cortisol, or 1 mg overnight dexamethasone suppression test) followed by determination of the cause (plasma ACTH, high-dose dexamethasone suppression test, CRH stimulation test, MRI pituitary, CT adrenals, inferior petrosal sinus sampling).

\begin{itemize}[noitemsep]
  \item Comprehensive medical history and physical examination
  \item Laboratory tests to assess organ function
  \item Imaging studies to visualize affected structures
  \item Specialized testing depending on the affected system
  \item Consultation with relevant specialists
\end{itemize}
\section{Prevention}
\begin{itemize}[noitemsep]
  \item Maintain a healthy lifestyle with balanced diet and exercise
  \item Avoid known risk factors
  \item Attend regular medical check-ups for early detection
  \item Follow recommended screening guidelines
\end{itemize}
\section{Treatment Options}
The right treatment depends on the cause: transsphenoidal surgery for Cushing's disease, surgical removal of adrenal tumors, chemotherapy for adrenal carcinoma, and ketoconazole or metyrapone for medical management. Symptomatic management and supportive care Enzyme replacement therapy for certain conditions Gene therapy (emerging for select conditions) Dietary modifications (e.g., phenylketonuria) Regular monitoring for associated complications Physical and occupational therapy Genetic counseling for family planning Participation in clinical trials for new therapies

\begin{itemize}[noitemsep]
  \item Medications to manage symptoms and address underlying mechanisms
  \item Lifestyle modifications and self-care measures
  \item Physical therapy and rehabilitation
  \item Surgical intervention when indicated
  \item Regular monitoring and follow-up care
\end{itemize}
\section{Living with It}
\begin{itemize}[noitemsep]
  \item Follow your treatment plan as prescribed
  \item Maintain regular follow-up with your healthcare provider
  \item Make necessary lifestyle adjustments
  \item Seek support from family, friends, or support groups
  \item Stay informed about your condition
\end{itemize}
\section{Outlook}
The outlook is generally positive with successful treatment.
